% !TeX root = ../VCL note.tex
\section{显示（Display）}

\subsection{显示器的种类和发展历史}

\begin{enumerate}
    \item 矢量显示（阴极射线管）
    \item 二维显示（电子束）
    \begin{itemize}
        \item 像素：每个可以点亮的屏幕点
        \item 分辨率：品目像素格式
        \item 纵横比：像素列数除以行数
        \item 帧：显示的一个画面
        \item 隔行扫描：保证亮度均衡
    \end{itemize}
    \item 彩色二维显示（彩色荧光幕，不同荧光剂分区涂布）
    \item 二维显示（LCD）：统一发光，用液晶部分控制光的偏振，调节亮度
    \item 二维显示（OLED）：有机发光半导体显示器。独立的LED单元。
\end{enumerate}

\subsection{色域}

美国国家电视标准委员会的色域标准称为NTSC。

\begin{itemize}
    \item sRGB：微软，HP（72% NSTC）
    \item ARGB: Adobe颜色（设计软件）
    \item Rec.2020:4K(电视的标准)
    \item DCI-P3 : Apple标准
\end{itemize}

\subsection{像素标准}

显卡的存储和刷新越厉害，达到4K8K就越容易。

\subsection{立体显示}

立体显示依据的原理是双目视觉。所有的立体显示技术基于人类的生理特征。最开始的显示技术是双目偏振片的3D显示。沉浸式的3D显示技术（AR,VR,XR）跟踪头显，裸眼3D+眼睛跟踪。现有的技术向MR-mix方向进行发展。

裸眼3D技术：LCD屏幕的方法：微柱透镜显示技术。利用视差屏障，令OLED屏幕上不同的像素点用不同角度进入人眼。

全息显示技术（Hplographic display）：用干涉记录波前的信息（数字全息摄影），用衍射记录物光波场（受限于较小的范围）。

多模态显示技术：触觉显示（触觉输出），嗅觉显示（控制挥发的成分。。。？）

